\sgaStableTarget{sga6.e06.sec.4}\section*{4. Computation of the @@M0@@ of a flag bundle}

\sgaStableTarget{sga6.e06.sec.4.1}4.1.
Let @@M1@@ be a scheme, and let @@M2@@ be an @@M3@@-Module locally free of rank @@M4@@. Let @@M5@@ be a sequence of integers @@M6@@ such that:
@@M7@@
Recall that the flag functor of type @@M8@@ of @@M9@@ is the functor on the category of @@M10@@-schemes, with values in the category of sets, which associates to an @@M11@@-scheme @@M12@@ the set of sequences of quotients of @@M13@@:
@@M14@@
where, for every @@M15@@, @@M16@@ is a quotient of @@M17@@, locally free of rank @@M18@@.

Such a functor is representable by an @@M19@@-scheme @@M20@@, endowed with a canonical sequence of quotients of the inverse image of @@M21@@, locally free of ranks @@M22@@

If @@M23@@, then @@M24@@.

If @@M25@@, then @@M26@@, the Grassmannian of type @@M27@@ of @@M28@@.

If all the @@M29@@ are equal to @@M30@@, one obtains the usual flag bundle of @@M31@@, denoted by @@M32@@.

\sgaStableTarget{sga6.e06.lem.4-2}\textbf{Lemma 4.2.} \uline{Let} @@M33@@ \uline{be a scheme}, @@M34@@ \uline{an} @@M35@@\uline{-Module locally free of rank} @@M36@@, @@M37@@, \uline{with} @@M38@@; \uline{let} @@M39@@ \uline{be two integers} @@M40@@ \uline{such that} @@M41@@; \uline{put} @@M42@@, @@M43@@. \uline{Let} @@M44@@ \uline{be the canonical quotients of the inverse image of} @@M45@@ \uline{on} @@M46@@, \uline{and let} @@M47@@ \uline{be the kernel of} @@M48@@. \uline{Then there is a canonical} @@M49@@-\uline{isomorphism}:
